Se consideraron medios homogéneos, isótropos y no magnéticos a 3 GHz. Los parámetros de permitividad y conductividad se tomaron del programa original y se trataron como entradas del modelo; no se midieron muestras físicas ni se cuantificó la variabilidad de dichas propiedades.

La atenuación se estimó en ventanas radiales declaradas. El balance se evaluó en un anillo sin fuentes, separado del conductor y de las fronteras absorbentes. Las comparaciones describieron las configuraciones simuladas, con la longitud del monopolo y el dominio escalados por la longitud de onda de cada medio.

No se resolvieron interfaces entre dos medios de propagación, heterogeneidad interna, anisotropía, calentamiento, SAR calibrada, impedancia de entrada, patrón de radiación completo ni eficiencia global. Tampoco se efectuó una validación experimental o contra datos independientes de enlaces. Estas limitaciones se retomaron al evaluar los objetivos en la discusión.
